
\documentclass{article}
\usepackage{colortbl}
\usepackage{makecell}
\usepackage{multirow}
\usepackage{supertabular}

\begin{document}

\newcounter{utterance}

\centering \large Interaction Transcript for game `dond', experiment `semi', episode 3 with claude{-}opus{-}4{-}8{-}azure.
\vspace{24pt}

{ \footnotesize  \setcounter{utterance}{1}
\setlength{\tabcolsep}{0pt}
\begin{supertabular}{c@{$\;$}|p{.15\linewidth}@{}p{.15\linewidth}p{.15\linewidth}p{.15\linewidth}p{.15\linewidth}p{.15\linewidth}}
    \# & $\;$A & \multicolumn{4}{c}{Game Master} & $\;\:$B\\
    \hline

    \theutterance \stepcounter{utterance}  
    & & \multicolumn{4}{p{0.6\linewidth}}{
        \cellcolor[rgb]{0.9,0.9,0.9}{
            \makecell[{{p{\linewidth}}}]{
                \texttt{\tiny{[P1$\langle$GM]}}
                \texttt{Igrate kooperativnu pregovaračku igru u kojoj treba da se dogovorite kako da podelite skup predmeta između vas i drugog igrača.} \\
\\ 
\texttt{Pravila:} \\
\texttt{(a) Vama i drugom igraču dodeljuje se skup predmeta. Svako od vas takođe dobija tajnu funkciju vrednosti, koja predstavlja koliko vrednujete svaku vrstu predmeta.} \\
\texttt{(b) Razmenjujete poruke sa drugim igračem kako biste se dogovorili ko dobija koje predmete. Svako od vas može poslati najviše 5 poruka, ili možete ranije završiti tako što ćete u bilo kom trenutku dati tajnu ponudu.} \\
\texttt{(c) Od svakog od vas se traži da podnese tajnu ponudu koja navodi predmete koje želite, formatiranu u uglastim zagradama na sledeći način: "{[}Predlog: <broj> <naziv predmeta>, <broj> <naziv predmeta>, <...>{]}"} \\
\texttt{(d) Ako su vaše ponude komplementarne, tj. ima dovoljno predmeta da se ispune obe ponude, svakom igraču se dodeljuje rezultat na osnovu zbira vrednosti predmeta koje je dobio. U suprotnom, oboje dobijate nula poena.} \\
\texttt{(e) Vaš cilj je da maksimizujete zbir vašeg rezultata i rezultata drugog igrača.} \\
\\ 
\texttt{Počnimo.} \\
\\ 
\texttt{Skup dostupnih predmeta je:} \\
\\ 
\texttt{1 šargarepa, 4 torbe, 1 luk} \\
\\ 
\texttt{Vaša tajna funkcija vrednosti je:} \\
\\ 
\texttt{šargarepa: 3, torba: 1, luk: 3} \\
\\ 
\texttt{VAŽNO: Vaše poruke, osim ako je u pitanju tajna ponuda, direktno se prosleđuju drugom igraču, zato nemojte uključivati nikakav odgovor na pravila niti tekst koji najavljuje vašu poruku. Da biste dali tajnu ponudu, koristite naznačeni format. Ne koristite uglaste zagrade kada komunicirate sa drugim igračem ili će to biti protumačeno kao vaša tajna ponuda.} \\
            }
        }
    }
    & \\ \\

    \theutterance \stepcounter{utterance}  
    & \multicolumn{4}{p{0.6\linewidth}}{
        \cellcolor[rgb]{0.9,0.9,0.9}{
            \makecell[{{p{\linewidth}}}]{
                \texttt{\tiny{[P1$\rangle$GM]}}
                \texttt{Zdravo! Hajde da podelimo predmete tako da obojica budemo zadovoljni. Meni su najvažnije šargarepa i luk, dok mi torbe manje znače. Šta je tebi najvrednije? Ako tebi torbe više odgovaraju, rado bih uzeo šargarepu i luk, a tebi prepustio torbe.} \\
            }
        }
    }
    & & \\ \\

    \theutterance \stepcounter{utterance}  
    & & & \multicolumn{2}{p{0.3\linewidth}}{
        \cellcolor[rgb]{0.9,0.9,0.9}{
            \makecell[{{p{\linewidth}}}]{
                \texttt{\tiny{[GM$|$GM]}}
                \texttt{continue} \\
            }
        }
    }
    & & \\ \\

    \theutterance \stepcounter{utterance}  
    & & & \multicolumn{2}{p{0.3\linewidth}}{
        \cellcolor[rgb]{0.9,0.9,0.9}{
            \makecell[{{p{\linewidth}}}]{
                \texttt{\tiny{[GM$|$GM]}}
                \texttt{Zdravo! Hajde da podelimo predmete tako da obojica budemo zadovoljni. Meni su najvažnije šargarepa i luk, dok mi torbe manje znače. Šta je tebi najvrednije? Ako tebi torbe više odgovaraju, rado bih uzeo šargarepu i luk, a tebi prepustio torbe.} \\
            }
        }
    }
    & & \\ \\

    \theutterance \stepcounter{utterance}  
    & & \multicolumn{4}{p{0.6\linewidth}}{
        \cellcolor[rgb]{0.9,0.9,0.9}{
            \makecell[{{p{\linewidth}}}]{
                \texttt{\tiny{[GM$\rangle$P2]}}
                \texttt{Igrate kooperativnu pregovaračku igru u kojoj treba da se dogovorite kako da podelite skup predmeta između vas i drugog igrača.} \\
\\ 
\texttt{Pravila:} \\
\texttt{(a) Vama i drugom igraču dodeljuje se skup predmeta. Svako od vas takođe dobija tajnu funkciju vrednosti, koja predstavlja koliko vrednujete svaku vrstu predmeta.} \\
\texttt{(b) Razmenjujete poruke sa drugim igračem kako biste se dogovorili ko dobija koje predmete. Svako od vas može poslati najviše 5 poruka, ili možete ranije završiti tako što ćete u bilo kom trenutku dati tajnu ponudu.} \\
\texttt{(c) Od svakog od vas se traži da podnese tajnu ponudu koja navodi predmete koje želite, formatiranu u uglastim zagradama na sledeći način: "{[}Predlog: <broj> <naziv predmeta>, <broj> <naziv predmeta>, <...>{]}"} \\
\texttt{(d) Ako su vaše ponude komplementarne, tj. ima dovoljno predmeta da se ispune obe ponude, svakom igraču se dodeljuje rezultat na osnovu zbira vrednosti predmeta koje je dobio. U suprotnom, oboje dobijate nula poena.} \\
\texttt{(e) Vaš cilj je da maksimizujete zbir vašeg rezultata i rezultata drugog igrača.} \\
\\ 
\texttt{Počnimo.} \\
\\ 
\texttt{Skup dostupnih predmeta je:} \\
\\ 
\texttt{1 šargarepa, 4 torbe, 1 luk} \\
\\ 
\texttt{Vaša tajna funkcija vrednosti je:} \\
\\ 
\texttt{šargarepa: 0, torba: 1, luk: 6} \\
\\ 
\texttt{VAŽNO: Vaše poruke, osim ako je u pitanju tajna ponuda, direktno se prosleđuju drugom igraču, zato nemojte uključivati nikakav odgovor na pravila niti tekst koji najavljuje vašu poruku. Da biste dali tajnu ponudu, koristite naznačeni format. Ne koristite uglaste zagrade kada komunicirate sa drugim igračem ili će to biti protumačeno kao vaša tajna ponuda.} \\
\\ 
\texttt{Zdravo! Hajde da podelimo predmete tako da obojica budemo zadovoljni. Meni su najvažnije šargarepa i luk, dok mi torbe manje znače. Šta je tebi najvrednije? Ako tebi torbe više odgovaraju, rado bih uzeo šargarepu i luk, a tebi prepustio torbe.} \\
            }
        }
    }
    & \\ \\

    \theutterance \stepcounter{utterance}  
    & & & \multicolumn{4}{p{0.6\linewidth}}{
        \cellcolor[rgb]{0.9,0.9,0.9}{
            \makecell[{{p{\linewidth}}}]{
                \texttt{\tiny{[GM$\langle$P2]}}
                \texttt{Zvuči kao dobar početak. Slažem se da tebi pripadnu šargarepa. Ali moram da budem iskren — meni je luk zapravo dosta vredan, dok mi torbe donose malu vrednost. Da bismo zajedno izvukli najviše, reci mi koliko ti tačno vredn} \\
            }
        }
    }
     \\ \\

    \theutterance \stepcounter{utterance}  
    & & & \multicolumn{2}{p{0.3\linewidth}}{
        \cellcolor[rgb]{0.9,0.9,0.9}{
            \makecell[{{p{\linewidth}}}]{
                \texttt{\tiny{[GM$|$GM]}}
                \texttt{continue} \\
            }
        }
    }
    & & \\ \\

    \theutterance \stepcounter{utterance}  
    & & & \multicolumn{2}{p{0.3\linewidth}}{
        \cellcolor[rgb]{0.9,0.9,0.9}{
            \makecell[{{p{\linewidth}}}]{
                \texttt{\tiny{[GM$|$GM]}}
                \texttt{Zvuči kao dobar početak. Slažem se da tebi pripadnu šargarepa. Ali moram da budem iskren — meni je luk zapravo dosta vredan, dok mi torbe donose malu vrednost. Da bismo zajedno izvukli najviše, reci mi koliko ti tačno vredn} \\
            }
        }
    }
    & & \\ \\

    \theutterance \stepcounter{utterance}  
    & & \multicolumn{4}{p{0.6\linewidth}}{
        \cellcolor[rgb]{0.9,0.9,0.9}{
            \makecell[{{p{\linewidth}}}]{
                \texttt{\tiny{[P1$\langle$GM]}}
                \texttt{Zvuči kao dobar početak. Slažem se da tebi pripadnu šargarepa. Ali moram da budem iskren — meni je luk zapravo dosta vredan, dok mi torbe donose malu vrednost. Da bismo zajedno izvukli najviše, reci mi koliko ti tačno vredn} \\
            }
        }
    }
    & \\ \\

    \theutterance \stepcounter{utterance}  
    & \multicolumn{4}{p{0.6\linewidth}}{
        \cellcolor[rgb]{0.9,0.9,0.9}{
            \makecell[{{p{\linewidth}}}]{
                \texttt{\tiny{[P1$\rangle$GM]}}
                \texttt{Razumem — ako je tebi luk vredniji, ima smisla da ga ti uzmeš. Meni luk vredi 3, ali ako je tebi više, neka bude tvoj. S obzirom da tebi torbe donose malu vrednost, predlažem da ja uzmem šargarepu i sve četiri torbe, a tebi ostane luk. Tako svako dobija ono što najviše ceni i zajedno izvlačimo maksimum. Da li ti to odgovara?} \\
            }
        }
    }
    & & \\ \\

    \theutterance \stepcounter{utterance}  
    & & & \multicolumn{2}{p{0.3\linewidth}}{
        \cellcolor[rgb]{0.9,0.9,0.9}{
            \makecell[{{p{\linewidth}}}]{
                \texttt{\tiny{[GM$|$GM]}}
                \texttt{continue} \\
            }
        }
    }
    & & \\ \\

    \theutterance \stepcounter{utterance}  
    & & & \multicolumn{2}{p{0.3\linewidth}}{
        \cellcolor[rgb]{0.9,0.9,0.9}{
            \makecell[{{p{\linewidth}}}]{
                \texttt{\tiny{[GM$|$GM]}}
                \texttt{Razumem — ako je tebi luk vredniji, ima smisla da ga ti uzmeš. Meni luk vredi 3, ali ako je tebi više, neka bude tvoj. S obzirom da tebi torbe donose malu vrednost, predlažem da ja uzmem šargarepu i sve četiri torbe, a tebi ostane luk. Tako svako dobija ono što najviše ceni i zajedno izvlačimo maksimum. Da li ti to odgovara?} \\
            }
        }
    }
    & & \\ \\

    \theutterance \stepcounter{utterance}  
    & & \multicolumn{4}{p{0.6\linewidth}}{
        \cellcolor[rgb]{0.9,0.9,0.9}{
            \makecell[{{p{\linewidth}}}]{
                \texttt{\tiny{[GM$\rangle$P2]}}
                \texttt{Razumem — ako je tebi luk vredniji, ima smisla da ga ti uzmeš. Meni luk vredi 3, ali ako je tebi više, neka bude tvoj. S obzirom da tebi torbe donose malu vrednost, predlažem da ja uzmem šargarepu i sve četiri torbe, a tebi ostane luk. Tako svako dobija ono što najviše ceni i zajedno izvlačimo maksimum. Da li ti to odgovara?} \\
            }
        }
    }
    & \\ \\

    \theutterance \stepcounter{utterance}  
    & & & \multicolumn{4}{p{0.6\linewidth}}{
        \cellcolor[rgb]{0.9,0.9,0.9}{
            \makecell[{{p{\linewidth}}}]{
                \texttt{\tiny{[GM$\langle$P2]}}
                \texttt{Odlično, to mi savršeno odgovara. Ti uzimaš šargarepu i sve četiri torbe, a ja luk. Tako svako dobija ono što najviše ceni. Dogovoreno!} \\
\\ 
\texttt{{[}Predlog: 1 luk{]}} \\
            }
        }
    }
     \\ \\

    \theutterance \stepcounter{utterance}  
    & & & \multicolumn{2}{p{0.3\linewidth}}{
        \cellcolor[rgb]{0.9,0.9,0.9}{
            \makecell[{{p{\linewidth}}}]{
                \texttt{\tiny{[GM$|$GM]}}
                \texttt{proposal} \\
            }
        }
    }
    & & \\ \\

    \theutterance \stepcounter{utterance}  
    & & & \multicolumn{2}{p{0.3\linewidth}}{
        \cellcolor[rgb]{0.9,0.9,0.9}{
            \makecell[{{p{\linewidth}}}]{
                \texttt{\tiny{[GM$|$GM]}}
                [0, 0, 1]
            }
        }
    }
    & & \\ \\

    \theutterance \stepcounter{utterance}  
    & & \multicolumn{4}{p{0.6\linewidth}}{
        \cellcolor[rgb]{0.9,0.9,0.9}{
            \makecell[{{p{\linewidth}}}]{
                \texttt{\tiny{[P1$\langle$GM]}}
                \texttt{Drugi igrač je podneo tajnu ponudu. Sada je na vama red da podnesete sopstvenu ponudu.} \\
            }
        }
    }
    & \\ \\

    \theutterance \stepcounter{utterance}  
    & \multicolumn{4}{p{0.6\linewidth}}{
        \cellcolor[rgb]{0.9,0.9,0.9}{
            \makecell[{{p{\linewidth}}}]{
                \texttt{\tiny{[P1$\rangle$GM]}}
                \texttt{{[}Predlog: 1 šargarepa, 4 torbe{]}} \\
            }
        }
    }
    & & \\ \\

    \theutterance \stepcounter{utterance}  
    & & & \multicolumn{2}{p{0.3\linewidth}}{
        \cellcolor[rgb]{0.9,0.9,0.9}{
            \makecell[{{p{\linewidth}}}]{
                \texttt{\tiny{[GM$|$GM]}}
                \texttt{proposal} \\
            }
        }
    }
    & & \\ \\

    \theutterance \stepcounter{utterance}  
    & & & \multicolumn{2}{p{0.3\linewidth}}{
        \cellcolor[rgb]{0.9,0.9,0.9}{
            \makecell[{{p{\linewidth}}}]{
                \texttt{\tiny{[GM$|$GM]}}
                [1, 4, 0]
            }
        }
    }
    & & \\ \\

    \theutterance \stepcounter{utterance}  
    & & & \multicolumn{2}{p{0.3\linewidth}}{
        \cellcolor[rgb]{0.9,0.9,0.9}{
            \makecell[{{p{\linewidth}}}]{
                \texttt{\tiny{[GM$|$GM]}}
                [[1, 4, 0], [0, 0, 1]]
            }
        }
    }
    & & \\ \\

\end{supertabular}
}

\end{document}
